\chapter{실물 제작}\label{ch:hardware}

이 장이 프로젝트의 최종 목표다. 여기서 만드는 두 기계는 앞 장의 FPGA 구현과 함께 세 대의 전시물을 이룬다.
기계식, 릴레이식, 전자식의 순서로 놓으면 같은 계산을 하는 기계가 어떻게 발전했는지를 보여 준다.
소프트웨어 모델은 세 대 모두의 기준이고, FPGA는 전시물인 동시에 두 실물의 구동 제어와 결과 대조도 맡는다.

실물은 같은 봄베를 두 가지 방식으로 만든다. 실제 영국 봄베처럼 드럼이 물리적으로 회전하는 \emph{기어식}과,
미 육군 봄베처럼 회전 부품 없이 릴레이로 구성하는 \emph{릴레이식}이다. 이 장은 제작이 진행된 뒤에 쓴다.
여기에는 이미 정해진 설계 방침과 설계 요구사항만 적는다.

\section{설계 요구사항: 전시와 실습}

실물은 전시관에 두고 관람객이 직접 조작해 볼 수 있어야 한다. 전시 여부는 아직 정해지지 않았지만 이를
설계 요구사항으로 둔다(표~\ref{tab:exhibit}). 세부는 설계 단계에서 확정한다.

\begin{table}[ht]
\centering
\caption{전시와 실습을 위한 설계 요구사항.}
\label{tab:exhibit}
\small
\begin{tabularx}{\linewidth}{>{\raggedright\arraybackslash}p{26mm}>{\raggedright\arraybackslash}X}
\toprule
영역 & 요구 \\
\midrule
실습 흐름 & 관람객이 실제 운용자의 순서를 따른다. 준비된 암호문과 크립 카드를 고르고, 패치 케이블로 메뉴를
  배선하고, 구동하고, 정지를 확인하고, 체킹한다. 한 번의 실습이 전시 동선에 맞는 시간 안에 끝나도록 축소한
  시나리오(짧은 메뉴, 제한된 위치 범위)를 준비한다 \\
안전 & 인체에 안전한 저전압만 노출한다. 회전부와 기어에는 투명 덮개와 인터록(덮개를 열면 정지)을 두고,
  비상 정지 버튼을 단다. 끼임, 화상, 감전의 위험 요소를 없앤다 \\
내구성과 유지보수 & 장시간의 반복 구동을 견디는 접점, 브러시, 기어 재질을 쓴다. 소모품의 교체가 쉬운 구조로
  하고, 기준 메뉴로 알려진 정지 위치를 확인하는 자가 진단 모드를 둔다 \\
가시성과 설명 & 드럼의 회전과 정지가 눈에 보이도록 배치하고, 전류가 흐르는 경로를 램프 등으로 표시하며,
  설명 패널을 둔다. 해설은 이 보고서와 입문 가이드의 원리 설명을 다시 쓴다 \\
복원 & 실습 뒤 배선과 드럼을 초기 상태로 쉽게 되돌린다 \\
전자 장치의 역할 & 소프트웨어와 FPGA는 실물 뒤에서 기준 모델(실물 결과의 대조), 정지 판정의 보조, 진단,
  관람객용 표시(해독 결과 화면)를 맡는다. 기계식과 릴레이식에서 해독 동작 자체는 기계가 한다 \\
세 기계의 비교 & 같은 메뉴 카드를 기계식, 릴레이식, FPGA에 걸어 같은 정지가 나오는 것을 보인다. 위치 하나를
  시험하는 데 걸리는 시간의 차이를 실측해 표시한다. FPGA는 움직이는 부품이 없으므로 내부 상태(676가닥의 선,
  드럼 위치, 정지)를 LED나 화면으로 보여 주는 장치를 둔다 \\
\bottomrule
\end{tabularx}
\end{table}

세 방식을 모두 유지한다. 관람객에게 실제로 돌아가는 기계가 보여야 하므로 기어식이 중심이고, 릴레이식과
FPGA는 같은 원리가 다른 기술로 옮겨 간 모습을 보여 준다.

\section{두 방식의 공통 구조}

두 방식 모두 각 위치에서 가설이 모순인지 판정하는 것은 배선을 흐르는 전류다. 차이는 위치를 어떻게
진행시키고 치환을 어떻게 바꾸는가에 있다(표~\ref{tab:tracks}). 그래서 두 방식의 스크램블러가 같은
인터페이스(26선 입력, 26선 출력, 위치 진행 신호)를 갖게 하고, 메뉴 배선과 대각선 보드와 정지 감지는
한 번만 만들어 함께 쓴다. 이렇게 하면 같은 메뉴를 두 방식으로 돌려 결과를 직접 비교할 수 있다.

\begin{table}[ht]
\centering
\caption{기어식과 릴레이식의 역할 분담.}
\label{tab:tracks}
\begin{tabular}{lll}
\toprule
기능 & 기어식 & 릴레이식 \\
\midrule
위치 진행 & 모터, 기어열, 단 사이 자리올림 & 릴레이 계수기 (26진 세 자리) \\
스크램블링 & 회전 드럼의 브러시와 접점판 & 릴레이 접점망 \\
메뉴 배선, 대각선 보드 & 공용 & 공용 \\
정지 감지 & 공용 (릴레이) & 공용 (릴레이) \\
\bottomrule
\end{tabular}
\end{table}

릴레이식의 근거는 미 육군 봄베다. 이 기계는 로터의 내부 배선을 모사한 고정 배선 회로를 두고, 전류를
그 회로의 어느 접점으로 들여보낼지를 릴레이로 정했다. 위치 진행은 릴레이 세 쌍의 조합으로 만드는 순환
시프트였다고 전한다~\cite{dturing2021army}. 로터 한 단의 치환이 ``시프트, 고정 배선, 역방향 시프트''로
나뉜다는 것은 식~\eqref{eq:through}의 꼴 그대로다.

실물 단계에서 FPGA는 논리를 대신하지 않는다. 스크램블링은 드럼 또는 릴레이가 직접 하고, FPGA는 구동
제어와 기준 모델과의 결과 대조에 쓴다.

\begin{todobox}{공통부}
스크램블러 인터페이스의 정의, 정지 감지 회로, 메뉴 배선 패널의 설계와 시험 결과.
\end{todobox}

\begin{todobox}{기어식}
드럼과 접점판의 시제품(재질, 치수, 브러시 접촉의 안정성), 드럼 세 개와 기어 자리올림으로 이루어진
스크램블러 한 유닛, 17,576 위치 순회와 치환의 대조 결과. 실제 봄베의 드럼 구조와 구동 기구에 관한
자료 가운데 확인하지 못한 항목(브러시와 접점의 재질과 치수, 기어비, 스크램블러에서 케이블과 반사판이
붙는 쪽)의 조사 결과도 여기에 적는다.
\end{todobox}

\begin{todobox}{릴레이식}
로터 한 단의 릴레이 접점망(시프트를 몇 단으로, 어떤 단위로 구성하는지와 접점 수), 위치 계수기,
스크램블러 한 유닛의 대조 결과.
\end{todobox}

\begin{todobox}{전시와 실습의 설계}
표~\ref{tab:exhibit}의 요구사항을 구체화한 결과. 축소 시나리오(메뉴와 위치 범위, 한 회 소요 시간),
안전 장치와 인터록, 전류 경로의 표시 방법, 자가 진단 절차, 초기화 절차, 실물과 기준 모델을 대조하는
전자 장치의 구성.
\end{todobox}

\begin{todobox}{세 방식의 비교}
같은 메뉴를 기계식, 릴레이식, FPGA로 돌린 정지 목록의 일치와 위치당 시간, 부품 수, 신뢰성의 비교.
\end{todobox}
